\section{A Runtime Control Plane for Tabular QA}
\label{sec:architecture}

\subsection{Control-Plane Architecture}

MOMO defines a control plane with three roles. A \textbf{planner} orchestrates the task at a high level: it decomposes the question, decides what to search for, decides which sources to commit to, and produces the final answer. Two classes of \textbf{contract executors} do the data-touching work: a \textbf{search subagent} discovers candidate datasets, and an \textbf{inspect subagent} extracts evidence from a chosen source family. Under MOMO's delegation mode, the planner's tool surface is reduced to four tools: planning utilities, a final-answer tool, \texttt{search\_subagent(contract)}, and \texttt{inspect\_subagent(contract)}. Raw search, peek, SQL, code-execution, and download tools are removed from the planner. This single restriction is what makes the planner an \emph{orchestrator}.

Both executor classes are spawned fresh per contract, run to completion or budget exhaustion, and emit a single compact return. Their internal scratchpads --- any schemas, query attempts, tracebacks, or partial results --- are discarded on return. The planner never sees them. This is the \emph{context firewall}, the central system-level property of the control plane.

\subsection{Source-Discovery Contracts}

The search subagent receives a \texttt{SearchContract} with: \texttt{contract\_id} for traceability and retry linking; \texttt{search\_goal}, a natural-language objective; \texttt{required\_source\_traits}, structured hints such as ``contains county FIPS codes'' or ``covers 2019--2021''; \texttt{budget\_calls}, a hard upper bound; \texttt{constraints}, exclusion criteria and tie-breakers; \texttt{known\_context}, datasets or facts the planner has already established; and \texttt{requested\_budget\_calls}, the planner's preferred budget, which the runtime may revise downward.

The search worker has access to retrieval and lightweight peek/profile tools, but not to query, code execution, or download. Its return schema is: \texttt{status} $\in$ \{\texttt{success}, \texttt{partial}, \texttt{failed}, \texttt{budget\_exhausted}\}; \texttt{candidates}, a small list of source-family handles with brief justifications; \texttt{search\_summary}; \texttt{missing\_coverage}, traits not satisfied by any candidate; \texttt{retry\_recommendation}, structured suggestions for a refined retry contract; and \texttt{failure\_reason} on \texttt{failed}.

\subsection{Source-Commitment Contracts}

The inspect subagent receives an \texttt{InspectContract} with: \texttt{contract\_id}; \texttt{objective}; \texttt{source\_family\_ids}, explicit identifiers for the dataset family the worker is allowed to touch; \texttt{required\_outputs}, an enumerated list of evidence pieces; \texttt{success\_criteria}; \texttt{budget\_calls}; \texttt{constraints} (e.g., size or runtime caps); \texttt{known\_context}; \texttt{retry\_of\_contract\_id}, non-empty when this is a retry; and \texttt{requested\_budget\_calls}.

The inspect worker can use the full tabular tool set --- peek, read, grep, parse XML, query files, download, and execute code --- but cannot see planner tools or invoke search. Its return schema includes \texttt{answer\_fragments} keyed by entries in \texttt{required\_outputs}; \texttt{missing\_outputs}, the entries it could not produce; \texttt{evidence} pointers (file paths, row ranges, query digests); \texttt{executor\_summary}; and \texttt{retry\_recommended} with an optional structured hint.

The \textbf{partial-success protocol} is load-bearing: a worker that finds 2019 and 2020 data but not 2021 returns \texttt{status="partial"} with \texttt{answer\_fragments=\{2019:\dots,2020:\dots\}} and \texttt{missing\_outputs=[2021]}. The planner can act on \texttt{missing\_outputs} directly --- retry with a different source family, accept the gap, or report it in the final answer --- without inspecting a raw traceback.

\subsection{Runtime Guards}

Two \texttt{SteeringHandler}-based guards act as execution-discipline operators in the control plane.

The \textbf{\texttt{\_FileReferenceGuard}} blocks any data-tool call whose arguments reference a dataset only by logical ID, with no accompanying \texttt{s3\_uri} or \texttt{dataset\_id+file\_path} pair. It also blocks references outside the data-lake bucket allowlist. Bookkeeping errors surface immediately as contract violations rather than ``file not found'' retries that drain the worker's budget.

The \textbf{\texttt{\_InspectSourceGuard}} restricts an inspect worker to the \texttt{source\_family\_ids} declared in its contract. Tool calls referencing other datasets are blocked at the runtime layer, before the call reaches the underlying tool. Source drift is prevented structurally; the worker cannot rationalize its way into a tangential dataset.

\subsection{Context Isolation and Provenance}

The context firewall is implemented through bounded workers plus compact return tools. Each subagent calls its \texttt{contract\_return} tool exactly once, with a payload conforming to the return schemas above. Schemas, row dumps, long tracebacks, and failed SQL attempts never enter the planner's conversation. A \texttt{\_SubagentToolLedger} attached to each worker records subagent start and end events, per-call tool budgets, the list of tools invoked, token usage, and cost --- producing \textbf{provenance over agent execution traces} that is available to evaluation and post-hoc analysis but not visible to the planner.

A planner that issues four inspect contracts sees four compact returns totaling a few hundred tokens of structured evidence and missing-output lists. The same task in a single-agent baseline exposes the planner to four schema dumps (potentially thousands of tokens), several failed-query tracebacks, and partial row dumps from peek calls.

\subsection{Optional Runtime Controls}

Three companion features are exposed as separable feature flags and can be combined with delegation. They give the control plane additional levers without changing its core shape; appendix~\ref{sec:appendix} carries the detailed prompt and tool specifications.

\paragraph{\texttt{results}.} Wraps \texttt{peek\_file} and \texttt{peek\_multiple} to attach a precomputed profile (schema columns, table kind, LLM-authored description, snippet, optional row counts, top rows, and per-column statistics). The wrapper improves data-environment perception, reducing unnecessary inspection and query calls.

\paragraph{\texttt{cot}.} Forces structured pre-tool and post-tool records via a \texttt{cot} tool. A \texttt{CoTSteerHandler} plugin cancels tool calls that violate the protocol, replacing free-text chain-of-thought compliance with runtime enforcement.

\paragraph{\texttt{sprint}.} Single-agent commitment control. In \texttt{cadence} mode, every $k$ counted tool calls the agent must call \texttt{sprint(kind="cadence",\dots)} before continuing. In \texttt{commitment} mode, the agent must declare a source-commitment contract before working on a source and reflect on budget exhaustion or source switch. The \texttt{sprint} tool writes a persistent \texttt{\#\# CURRENT SPRINT} section into the agent's system prompt.

\paragraph{Delegation vs.\ sprint.} The two are alternative control-plane strategies and are mutually exclusive in the flag validator. \texttt{sprint} forces a single agent to commit and reflect; \texttt{delegation} lifts the commitment boundary up one level. The two expose different trade-offs between accuracy, cost, interruptiveness, and context hygiene, not strict improvements over each other.
